\hypertarget{appendix-8e5214819615}{}
\section*{Class access and invariants}
\section{Class access and invariants}\label{class-access-and-invariants}

\textbf{Curriculum depth:} Developed · \textbf{First seen:} C007 · \textbf{Latest development:} C009

\subsection{Definition and mechanism}\label{definition-and-mechanism}

Constructors establish externally usable state before publication; failed ordinary construction destroys completed subobjects. Generated memberwise moves do not automatically preserve relations such as extent-to-buffer consistency.

struct and class differ in default member and base-class access: public versus private. Either supports methods, inheritance, constructors and hidden data. The record-versus-invariant distinction is design convention; an invariant is a condition each public operation must preserve.

\subsection{Prerequisites}\label{prerequisites}

\begin{itemize}
\tightlist
\item
  \hyperref[appendix-2e4d494c2d94]{Declarations and definitions} --- Class declarations establish members and bases whose access must be public or controlled.
\end{itemize}

\subsection{Used by}\label{used-by}

\begin{itemize}
\tightlist
\item
  \hyperref[appendix-f5593256bcaf]{Logical constness and mutable caches}
\item
  \hyperref[appendix-7f8ec83d0886]{Constructors and invariant establishment}
\end{itemize}

\subsection{Recall and next depth}\label{recall-and-next-depth}

Give an example, explain why the rule exists, and identify which guarantee would be missing if the prerequisite were ignored.

\textbf{Pending:} polymorphic invariants and subsystem-level failure policies.

\subsection{Source and navigation}\label{source-and-navigation}

\hyperref[appendix-fd090d982a15]{C009 development}

\href{Continuous\%20Notes/Computing\%20Foundations\%20-\%20From\%20Bits\%20to\%20C++\%20and\%20Systems\%20APIs\#Public\%20records\%20and\%20protected\%20invariants}{C007 teaching chapter} · \hyperref[appendix-fd090d982a15]{Source coverage} · \href{}{Dependency Map (02-Foundations-II, appendix-4ad47bdd269a)}
